% ============================================================================
% WORKBOOK — Section 8.6: Circle Graphs (Additional)
% State-specific (FL, MN, TX, VA)
% Practice-focused reinforcement for the study guide topic
% ============================================================================
\section{Circle Graphs}
\topicTitle{Circle Graphs}

\begin{quickReview}{Circle Graphs}
A \textbf{circle graph} (pie chart) shows how a whole is divided into parts. The full circle represents $100\%$ or $360°$.

\begin{itemize}[leftmargin=*]
    \item Each \textbf{sector} represents one category.
    \item To find a sector's angle: $\text{percent} \times 360°$.
    \item To find the percent from an angle: $\dfrac{\text{angle}}{360°} \times 100\%$.
    \item All sectors must add to $100\%$ (or $360°$).
    \item To find the number in a category: $\text{percent} \times \text{total}$.
\end{itemize}

\medskip
\textbf{Example:} $40\%$ of $200$ students prefer soccer.
Number $= 0.40 \times 200 = 80$ students. Angle $= 0.40 \times 360° = 144°$.
\end{quickReview}

% ── Warm-Up ──────────────────────────────────────────────────────────────────
\begin{practiceBox}[funGreen]{\faSun~Warm-Up}

\practiceHeader[funGreen]{Percent to Angle (and Back)}

\begin{multicols}{2}
\prob $25\%$ of a circle $=$ \answerBlank[2cm] degrees
\answerExplain{$90°$}{$0.25 \times 360° = 90°$. One-quarter of the circle is $90°$.}

\prob $50\%$ of a circle $=$ \answerBlank[2cm] degrees
\answerExplain{$180°$}{$0.50 \times 360° = 180°$. Half the circle is $180°$.}

\prob $10\%$ of a circle $=$ \answerBlank[2cm] degrees
\answerExplain{$36°$}{$0.10 \times 360° = 36°$.}

\prob A sector of $120°$ $=$ \answerBlank[2cm] percent
\answerExplain{$33.\overline{3}\%$}{$\frac{120}{360} \times 100 = 33.\overline{3}\%$. This is one-third of the circle.}

\prob A sector of $72°$ $=$ \answerBlank[2cm] percent
\answerExplain{$20\%$}{$\frac{72}{360} \times 100 = 20\%$.}

\prob $75\%$ of a circle $=$ \answerBlank[2cm] degrees
\answerExplain{$270°$}{$0.75 \times 360° = 270°$. Three-quarters of the circle is $270°$.}
\end{multicols}
\end{practiceBox}

% ── Practice A: Reading Circle Graphs ────────────────────────────────────────
\begin{practiceBox}{Reading a Circle Graph}

\practiceHeader{A survey of $300$ students asked about their favorite season:}

\begin{center}
\begin{tabular}{lcc}
    \textbf{Season} & \textbf{Percent} & \textbf{Number of students} \\
    \hline
    Summer  & $40\%$ & \rule{2cm}{0.4pt} \\
    Fall    & $25\%$ & \rule{2cm}{0.4pt} \\
    Spring  & $20\%$ & \rule{2cm}{0.4pt} \\
    Winter  & $15\%$ & \rule{2cm}{0.4pt} \\
\end{tabular}
\end{center}

\prob How many students chose Summer? \answerBlank[3cm]
\answerExplain{$120$ students}{$0.40 \times 300 = 120$. Multiply the percent (as a decimal) by the total number of students.}

\prob How many students chose Fall? \answerBlank[3cm]
\answerExplain{$75$ students}{$0.25 \times 300 = 75$.}

\prob How many students chose Spring? \answerBlank[3cm]
\answerExplain{$60$ students}{$0.20 \times 300 = 60$.}

\prob How many chose Winter? \answerBlank[3cm]
\answerExplain{$45$ students}{$0.15 \times 300 = 45$. Check: $120 + 75 + 60 + 45 = 300$. \checkmark}

\prob What is the angle for the Summer sector? \answerBlank[3cm]
\answerExplain{$144°$}{$0.40 \times 360° = 144°$.}

\prob What is the angle for the Winter sector? \answerBlank[3cm]
\answerExplain{$54°$}{$0.15 \times 360° = 54°$.}

\end{practiceBox}

% ── Practice B: Finding Missing Percents ─────────────────────────────────────
\begin{practiceBox}[funTeal]{Finding Missing Percents and Angles}

\prob A circle graph has three sectors: $35\%$, $40\%$, and $x\%$. Find $x$. \answerBlank[3cm]
\answerExplain{$25\%$}{All sectors must total $100\%$: $35 + 40 + x = 100$, so $x = 25\%$.}

\prob A circle graph has four sectors with angles $100°$, $80°$, $60°$, and $y°$. Find $y$. \answerBlank[3cm]
\answerExplain{$120°$}{All angles total $360°$: $100 + 80 + 60 + y = 360$, so $y = 120°$.}

\prob A class of $40$ students voted on a trip. Museum: $12$, Zoo: $16$, Park: $8$, Aquarium: $4$. Find each category's percent.
\answerExplain{Museum $30\%$, Zoo $40\%$, Park $20\%$, Aquarium $10\%$}{Museum: $\frac{12}{40} = 30\%$. Zoo: $\frac{16}{40} = 40\%$. Park: $\frac{8}{40} = 20\%$. Aquarium: $\frac{4}{40} = 10\%$. Check: $30 + 40 + 20 + 10 = 100\%$. \checkmark}

\prob In Problem 15, what angle represents the Zoo sector? \answerBlank[3cm]
\answerExplain{$144°$}{Zoo is $40\%$. Angle $= 0.40 \times 360° = 144°$.}

\end{practiceBox}

% ── Practice C: Converting Between Representations ───────────────────────────
\begin{practiceBox}[funOrange]{Percents, Fractions, and Decimals}

\practiceHeader[funOrange]{Express each sector as a fraction, decimal, and percent.}

\prob A sector is $\frac{1}{5}$ of the circle. Write as a percent and find the angle. \answerBlank[3cm]
\answerExplain{$20\%$; $72°$}{$\frac{1}{5} = 0.20 = 20\%$. Angle $= 0.20 \times 360° = 72°$.}

\prob A sector is $\frac{3}{8}$ of the circle. Write as a percent and find the angle. \answerBlank[3cm]
\answerExplain{$37.5\%$; $135°$}{$\frac{3}{8} = 0.375 = 37.5\%$. Angle $= 0.375 \times 360° = 135°$.}

\prob A sector has an angle of $45°$. Write as a fraction and a percent. \answerBlank[3cm]
\answerExplain{$\frac{1}{8}$; $12.5\%$}{$\frac{45}{360} = \frac{1}{8} = 0.125 = 12.5\%$.}

\prob A sector represents $0.6$ of the data. Find the percent and angle. \answerBlank[3cm]
\answerExplain{$60\%$; $216°$}{$0.6 = 60\%$. Angle $= 0.6 \times 360° = 216°$.}

\end{practiceBox}

% ── True or False ────────────────────────────────────────────────────────────
\begin{practiceBox}[funPurple]{True or False?}

\trueOrFalse{All sectors of a circle graph must add up to $100\%$.}
\answerExplain{True}{A circle graph represents the whole data set. All parts must total $100\%$ (or $360°$).}

\trueOrFalse{A sector with a $90°$ angle represents $25\%$ of the data.}
\answerExplain{True}{$\frac{90}{360} = 0.25 = 25\%$. A $90°$ angle is one-quarter of the full circle.}

\trueOrFalse{If one sector takes up exactly half the circle, its angle is $180°$.}
\answerExplain{True}{Half of $360° = 180°$. Half the circle represents $50\%$ of the data.}

\end{practiceBox}

% ── Word Problems ────────────────────────────────────────────────────────────
\begin{practiceBox}[funBlue]{\faGlobe~Word Problems}

\wordProblem{A family's monthly budget is $\$4{,}000$. Housing is $35\%$, food is $20\%$, transportation is $15\%$, savings is $10\%$, and other is $20\%$. How much is spent on housing?}{dollars}
\answerExplain{$\$1{,}400$}{$0.35 \times 4{,}000 = 1{,}400$. The housing sector represents $35\%$ of the total budget.}

\wordProblem{A company's $500$ employees were surveyed about commuting. Driving: $55\%$, public transit: $25\%$, biking: $12\%$, walking: $8\%$. How many more employees drive than take public transit?}{employees}
\answerExplain{$150$ employees}{Drive: $0.55 \times 500 = 275$. Transit: $0.25 \times 500 = 125$. Difference: $275 - 125 = 150$.}

\wordProblem{A circle graph shows $30\%$ of $600$ books in a library are fiction. If $10\%$ of the fiction books are mysteries, how many mystery books are there?}{books}
\answerExplain{$18$ mystery books}{Fiction: $0.30 \times 600 = 180$ books. Mysteries: $0.10 \times 180 = 18$ books.}

\end{practiceBox}

% ── Challenge ────────────────────────────────────────────────────────────────
\begin{challengeBox}
\prob A school's circle graph shows that $\frac{2}{5}$ of students play a sport, $\frac{1}{4}$ are in band, and the rest are in drama or art club (split equally). If there are $200$ students, how many are in each activity? Find the angle for each sector.
\answerExplain{Sports: $80$ ($144°$); Band: $50$ ($90°$); Drama: $35$ ($63°$); Art: $35$ ($63°$)}{Sports: $\frac{2}{5} \times 200 = 80$. Band: $\frac{1}{4} \times 200 = 50$. Remaining: $200 - 80 - 50 = 70$; split equally: $35$ each. Angles: Sports $\frac{2}{5} \times 360° = 144°$. Band $\frac{1}{4} \times 360° = 90°$. Drama and Art each $\frac{35}{200} \times 360° = 63°$. Check: $144 + 90 + 63 + 63 = 360°$. \checkmark}
\end{challengeBox}

\encouragement{Great job reading and creating circle graphs! You can turn any data into a visual story!}
